% Arbre de Diane : transfert d'électrons du cuivre aux ions argent (corrigé de l'exercice 3, chapitre 3)
\begin{tikzpicture}[x=1cm,y=1cm,
  atome/.style={circle,minimum size=8.5mm,inner sep=0pt,font=\otSemi\fontsize{7.5}{9}\selectfont}]
\node[atome,fill=orange!70!red!80!black,text=white] (cu) at (0,0) {Cu};
\draw[otInk!70,line width=0.7pt,-{Stealth[length=2mm]}] (0.55,0) -- (1.45,0);
\node[atome,fill=cyan!35,text=otInk] (cu2) at (2.0,0) {Cu$^{2+}$};
\node[atome,fill=otMuted!18,text=otInk] (ag1) at (5.0,0.55) {Ag$^{+}$};
\node[atome,fill=otMuted!18,text=otInk] (ag2) at (5.0,-0.55) {Ag$^{+}$};
\draw[otInk!70,line width=0.7pt,-{Stealth[length=2mm]}] (5.55,0) -- (6.45,0);
\node[atome,fill=otMuted!55,text=white] at (7.0,0.55) {Ag};
\node[atome,fill=otMuted!55,text=white] at (7.0,-0.55) {Ag};
\draw[otGreenDark,line width=1.1pt,-{Stealth[length=2.4mm]}] (0.25,0.5) to[out=45,in=150] node[midway,above,pcLegende,text=otGreenDark] {un électron} (4.55,0.75);
\draw[otGreenDark,line width=1.1pt,-{Stealth[length=2.4mm]}] (0.25,-0.5) to[out=-45,in=-150] node[midway,below,pcLegende,text=otGreenDark] {un électron} (4.55,-0.75);
\node[pcLegende,anchor=north,align=center] at (1.0,-1.75) {oxydation :\\$\ce{Cu -> Cu^{2+} + 2e-}$};
\node[pcLegende,anchor=north,align=center] at (6.0,-1.75) {réduction (deux fois) :\\$\ce{Ag+ + e- -> Ag}$};
\end{tikzpicture}
\par\smallskip\legende{Un atome de cuivre cède deux électrons : il faut deux ions argent pour les capter, d'où le nombre 2 devant $\ce{Ag+}$ et $\ce{Ag}$.}
